\documentclass[profile=standard]{ipara-handbook}
\title{B05｜线性方程与线性微分方程：一点加核空间}
\author{}
\date{}
\begin{document}
\maketitle

\begin{quote}
本册回答：\textbf{为什么线性方程组和线性微分方程都会出现“特解 + 齐次通解”？这不是两门课碰巧使用同一句口诀，而是同一个线性映射逆像结构。}
\end{quote}

\section{本册定位：研究线性映射的纤维，而不是某种求解技巧}
线代 Topic04 Own $Ax=b$ 的可达性、核空间和仿射解集；高数 Topic12 Own 线性 ODE 的求解、初值与解族。B05 只 Own：
\[
\boxed{T^{-1}(b)=x_0+\ker T}
\]
这一跨有限维/函数空间都成立的线性结构。

它告诉我们：
\begin{itemize}
\item 一个特解负责“命中目标输出” $b$；
\item 核空间负责“怎样改变输入而不改变输出”；
\item 全部解就是一个特解沿所有不可见方向的平移。
\end{itemize}

\section{先固定三个对象：像、核与仿射平移}
设 $T:V\to W$ 为线性映射。定义
\[
\ker T:=\{x\in V:T(x)=0\},
\qquad
\operatorname{Im}T:=\{T(x):x\in V\}.
\]
核空间记录“哪些输入变化对输出完全不可见”，像空间记录“哪些目标输出真正可达”。若 $U\subseteq V$ 是线性子空间，$x_0+U:=\{x_0+u:u\in U\}$ 称为 $U$ 的一个仿射平移；它通常不经过原点，因此一般不是线性子空间。

有了这三个定义，B05 的母结构才有明确对象：
\[
\boxed{T^{-1}(b)\neq\varnothing\Longrightarrow T^{-1}(b)=x_0+\ker T.}
\]

\section{母定理：非空线性逆像一定是核空间的仿射平移}
若 $x_0$ 满足
\[
T(x_0)=b,
\]
则任意另一个解 $x$ 满足
\[
T(x-x_0)=T(x)-T(x_0)=0,
\]
所以
\[
x-x_0\in\ker T.
\]
反过来，对任意 $z\in\ker T$，
\[
T(x_0+z)=b.
\]
因此
\[
\boxed{T^{-1}(b)=x_0+\ker T.}
\]

这条证明只用了线性性，所以它比任何具体消元法、常系数特征方程或待定系数法都更根本。

\section{先分两道问题：能不能到达，到了以后有多少自由度}
线性方程真正包含两个逻辑上不同的问题：
\[
\boxed{b\in\operatorname{Im}T?}
\]
决定是否存在解；若存在，
\[
\boxed{\ker T}
\]
决定解之间还能怎样自由移动。

这也是为什么“有解”和“唯一解”不能混在一起：唯一解还要求
\[
\ker T=\{0\}.
\]

\section{线性方程组：非齐次解集为什么不是子空间}
对
\[
Ax=b,
\]
若 $x_0$ 是特解，则
\[
\boxed{x=x_0+z,\qquad z\in\ker A.}
\]
当 $b\neq0$ 时，解集通常不含原点，因此不是线性子空间，而是核空间的仿射平移。

所以齐次方程与非齐次方程的关系应读成：
\[
\boxed{\text{齐次方程决定方向，非齐次目标决定平移到哪里。}}
\]

\section{线性微分方程：函数空间里完全同一件事}
设线性微分算子
\[
L[y]=a_n(x)y^{(n)}+\cdots+a_1(x)y'+a_0(x)y.
\]
若
\[
L[y_p]=f,
\]
则任意齐次解 $y_h\in\ker L$ 满足
\[
L[y_p+y_h]=f.
\]
反过来，任意两个非齐次解之差都在 $\ker L$ 中。因此
\[
\boxed{\{y:L[y]=f\}=y_p+\ker L.}
\]

这里没有把微分算子“等同成一个矩阵”；我们只是使用了两者共享的线性映射结构。

\section{为什么 $n$ 阶齐次 ODE 会出现 $n$ 个自由常数}
在系数连续且最高阶系数在所研究区间不为零等常见正则条件下，$n$ 阶线性齐次 ODE 的解空间维数为 $n$。因此可取一组基础解
\[
y_1,\dots,y_n,
\]
使
\[
\ker L=\operatorname{span}\{y_1,\dots,y_n\}.
\]
于是非齐次通解写成
\[
\boxed{y=y_p+C_1y_1+\cdots+C_ny_n.}
\]

“$n$ 个常数”不是格式规定，而是核空间维数的坐标表现。

\section{初值条件：不是重新求方程，而是在仿射解族中切一刀}
例如 $n$ 阶 ODE 给定
\[
y(x_0),y'(x_0),\dots,y^{(n-1)}(x_0)
\]
时，取值算子 $y\mapsto y^{(j)}(x_0)$ 都是线性泛函；指定它们等于给定常数，就是作用在解空间上的仿射约束。在相应存在唯一性条件满足时，这 $n$ 个条件独立地锁定原本 $n$ 个自由常数，最终留下一个解。

因此解题结构可以压成：
\[
\boxed{\text{先构造整个仿射解族，再用初值选择其中一个点。}}
\]

\section{母例：两个世界只换对象，不换证明}
线代：
\[
\begin{bmatrix}1&1\end{bmatrix}x=1.
\]
一个特解是
\[
x_0=(1,0)^T,
\]
核空间由
\[
(1,-1)^T
\]
张成，所以
\[
x=(1,0)^T+t(1,-1)^T.
\]

ODE：
\[
y'-y=e^x.
\]
一个特解是
\[
y_p=xe^x,
\]
齐次核为
\[
\ker L=\{Ce^x:C\in\mathbb R\},
\]
因此
\[
y=(x+C)e^x.
\]

两题的对象、维数和求特解方法完全不同，但“一个点 + 核空间”的证明一字不改。

\section{为什么叠加原理与一点加核空间其实是同一件事}
线性算子满足
\[
L[\alpha y_1+\beta y_2]
=\alpha L[y_1]+\beta L[y_2].
\]
所以齐次解对线性组合封闭，形成核空间；非齐次解则不对任意线性组合封闭，却保持“两个解之差为齐次解”。

这说明：
\begin{itemize}
\item \textbf{叠加原理}解释核空间为什么是线性空间；
\item \textbf{一点加核空间}解释非齐次解集为什么是仿射空间。
\end{itemize}

\section{B05 串起哪些章节}
\begin{tabular}{p{0.23\linewidth}p{0.30\linewidth}p{0.39\linewidth}}
\hline
Owner / 下游 & 提供对象 & B05 的翻译责任\\
\hline
线代 Topic03/04 & image、kernel、$Ax=b$ 解空间 & 把“可达 + 自由度”抽成一般线性逆像结构\\
高数 Topic12 & 线性 ODE、齐次解、特解、初值 & 把函数解族解释为微分算子的仿射逆像\\
线代 Topic05 / ODE Extension & 特征根、谱、矩阵指数 & 只作为线性系统更深表示的扩展，不进入数学一标量 ODE 主干\\
\hline
\end{tabular}

\section{反桥梁：哪些“特解 + 通解”不能乱迁移}
\begin{tabular}{p{0.20\linewidth}p{0.04\linewidth}p{0.20\linewidth}p{0.48\linewidth}}
\hline
A & & B & 真正区别\\
\hline
非齐次解集 & $\neq$ & 线性子空间 & 一般是 $x_0+\ker T$，不经过原点\\
一个特解 & $\neq$ & 全部解 & 仍必须加上完整 kernel\\
线性 ODE & $\neq$ & 一般非线性 ODE & 非线性方程解集通常不具备仿射平移结构\\
微分算子 & $\neq$ & 有限维矩阵 & 两者共享线性结构，但对象空间与分析性质不同\\
\hline
\end{tabular}

\section{调用信号与一页压缩}
看到“已知一个解求全部解、非齐次通解、比较两个解、初值裁剪、齐次解空间”时，先问：
\begin{enumerate}
\item 当前算子是否线性？
\item 目标输出 $b$ 或 $f$ 是否可达？
\item 已知哪个特解？
\item 哪些输入变化会被算子送到零？
\item 附加条件还要从仿射解族中裁掉多少自由方向？
\end{enumerate}

最终只需记住
\[
\boxed{\text{非空线性逆像 = 一个特解 + 核空间}.}
\]

\end{document}
